\documentclass[10pt,letterpaper]{article}
\usepackage[margin=0.7in,headheight=14pt]{geometry}
\usepackage[T1]{fontenc}
\usepackage{lmodern,amsmath,amssymb,booktabs,xcolor,fancyhdr,microtype}
\usepackage[colorlinks=true,urlcolor=teal]{hyperref}
\definecolor{navy}{HTML}{192B4D}
\definecolor{teal}{HTML}{007D80}
\pagestyle{fancy}\fancyhf{}
\fancyhead[L]{\small MF 731 | Homework 4}
\fancyhead[R]{\small 09: Scenario Analysis}
\fancyfoot[C]{\small\thepage}
\setlength{\parindent}{0pt}
\setlength{\parskip}{5pt}
\hypersetup{pdftitle={MF 731 - Scenario Analysis}}
\begin{document}
{\Large\bfseries\textcolor{navy}{Scenario Analysis of a Delta-Hedged Option}}

\textbf{(a) Option price.} Write $\tau=T-t$. Under $Q$, the stock drifts are $r$,
and the log returns are jointly normal. Their moments give
\[
E^Q[(S_T^{(i)})^2\mid S_t^{(i)}=s_i]=s_i^2e^{(2r+\sigma_i^2)\tau},\qquad
E^Q[S_T^{(1)}S_T^{(2)}\mid S_t=s]=s_1s_2e^{(2r+\rho\sigma_1\sigma_2)\tau}.
\]
Expanding the squared payoff and discounting therefore gives
\[
\boxed{V(t,s_1,s_2;\sigma,\rho)=e^{r\tau}
\left(s_1^2e^{\sigma_1^2\tau}-2s_1s_2e^{\rho\sigma_1\sigma_2\tau}
+s_2^2e^{\sigma_2^2\tau}\right).}
\]
The deltas used for the initial hedge are
\[
\begin{aligned}
\partial_{s_1}V&=2e^{r\tau}\left(s_1e^{\sigma_1^2\tau}-s_2e^{\rho\sigma_1\sigma_2\tau}\right),\\
\partial_{s_2}V&=2e^{r\tau}\left(s_2e^{\sigma_2^2\tau}-s_1e^{\rho\sigma_1\sigma_2\tau}\right).
\end{aligned}
\]

\textbf{(b) Worst-case loss.} I use the supplied parameters, with
$r=0.0132$, $T=0.25$ and horizon $d=5/252$ years.
Initially,
\[
V(0)=921.739866,\qquad h_1=-21.952543,\qquad h_2=31.201877.
\]
The stock holdings are $(-h_1,-h_2)$, kept fixed for five days.
Using the portfolio definition in the question,
\[
\mathcal V_0=V(0)-h_1S_0^{(1)}-h_2S_0^{(2)}=-921.739866,
\]
\[
\mathcal V_d=V(d,S_0^{(1)}e^{x_1},S_0^{(2)}e^{x_2};\sigma',\rho')
-h_1S_0^{(1)}e^{x_1}-h_2S_0^{(2)}e^{x_2},\qquad
\ell=\mathcal V_0-\mathcal V_d.
\]
I enumerate all $6^4\times3=3{,}888$ scenarios and reprice with remaining
maturity $T-d$. The largest loss is $\boxed{1{,}000.36}$.

\begin{center}
\begin{tabular}{lrr}\toprule
Maximizing scenario & (b) Worst loss & (c) Worst weighted loss\\\midrule
Log returns $(x_1,x_2)$ & $(0.20,0.20)$ & $(0.05,0.05)$\\
Volatility factors $(\sigma_1'/\sigma_1,\sigma_2'/\sigma_2)$ & $(0.5,0.5)$ & $(0.75,0.75)$\\
New volatilities $(\sigma_1',\sigma_2')$ & $(0.1107,0.1518)$ & $(0.16605,0.2277)$\\
Correlation shift $\rho'-\rho$ & $+0.5$ & $+0.5$\\
New correlation $\rho'$ & $0.737$ & $0.737$\\
Portfolio loss $\ell$ & $1{,}000.36$ & $672.66$\\\bottomrule
\end{tabular}
\end{center}

\textbf{(c) Weighted scenario risk.} I multiply the five factor weights from
the question and maximize $w\ell$. The maximizing scenario is in the right
column above. Its weight is
\[
w=1\times1\times1.25\times1.25\times1.25=1.953125,
\qquad \max_n(w_n\ell_n)=\boxed{1{,}313.79}.
\]
Its actual portfolio loss is $\boxed{672.66}$.
The weights are used as given, without normalization. Since they can exceed
one, the weighted risk can exceed the worst unweighted loss.

\vfill
\textbf{Python code and running instructions:}\\
\url{https://github.com/jesuswalking/MF-731-HW1/tree/main/HW4}
\end{document}
